\chapter{Présentation de l'Organisme d'Accueil}

\section{Présentation d'OCP}
Fondé en 1920 sous le nom d'\textbf{Office Chérifien des Phosphates} et transformé en Société Anonyme (OCP S.A.) en 2008, le \textbf{Groupe OCP} est le leader mondial sur le marché de l'extraction, du traitement, de la transformation et de la commercialisation du phosphate et de ses dérivés (acide phosphorique, engrais minéraux phosphatés).

Le Maroc détient plus de 70\,\% des réserves prouvées de phosphate de la planète, ce qui confère au Groupe OCP un rôle géostratégique capital dans la sécurité alimentaire mondiale. Présent sur les cinq continents, le groupe emploie plus de 20\,000 collaborateurs et réalise un chiffre d'affaires annuel de plusieurs milliards de dollars.

Dans le cadre de sa vision prospective, le Groupe OCP a engagé depuis plus d'une décennie un vaste programme d'investissement industriel visant à doubler ses capacités d'extraction et à tripler sa production d'engrais, tout en adoptant une démarche rigoureuse de développement durable. Cette stratégie s'articule autour de la neutralité carbone d'ici 2040, de l'utilisation exclusive d'eaux non conventionnelles (dessalement d'eau de mer et recyclage des eaux usées) et de la production d'énergie propre à travers des parcs éoliens et photovoltaïques.

\section{Activités Industrielles}
Le Groupe OCP opère sur l'ensemble de la chaîne de valeur du phosphate, de l'amont minier jusqu'à la chimie fine de spécialité :

\begin{enumerate}
    \item \textbf{L'Extraction Minière :} Réalisée principalement à ciel ouvert sur les gisements historiques de Khouribga, Benguérir, Youssoufia (bassin des Gantour) et Boucraâ (Laâyoune). L'abattage du minerai brut s'effectue par décapage et dragage à l'aide de draglines et de pelles mécaniques de grande puissance.
    \item \textbf{La Préparation et l'Enrichissement du Minerai :}
    \begin{itemize}
        \item \textit{Concassage \& Criblage :} Réduction dimensionnelle et séparation granulométrique des roches phosphatées.
        \item \textit{Lavage \& Séchage :} Élimination des stériles argileux et régulation du taux d'humidité.
        \item \textit{Flottation :} Procédé physico-chimique permettant d'augmenter significativement la teneur en pentoxyde de phosphore ($P_2O_5$) en séparant sélectivement les carbonates et silicates.
    \end{itemize}
    \item \textbf{Le Transport Minéralier Continu :} Acheminement du phosphate enrichi soit par voie ferroviaire, soit par le système révolutionnaire de \textbf{Slurry Pipeline} reliant Khouribga à la plateforme chimique de Jorf Lasfar sur 235\,km, transportant le minerai sous forme de pulpe hydro-minérale avec une réduction massive de l'empreinte carbone.
    \item \textbf{La Transformation Chimique :} Réalisée sur les méga-complexes industriels de Jorf Lasfar et Safi. Le phosphate y est attaqué à l'acide sulfurique pour produire de l'acide phosphorique ($H_3PO_4$), servant de base à la synthèse d'engrais solides complexes de renommée mondiale : DAP (Di-Ammonium Phosphate), MAP (Mono-Ammonium Phosphate) et TSP (Triple Super Phosphate).
\end{enumerate}

\section{Site d'Accueil}
Le stage d'immersion a été réalisé au sein du complexe industriel d'OCP à \textbf{[SITE INDUSTRIEL OCP]}. Ce pôle industriel intègre des installations complexes de traitement minier, de convoyage longue distance et de stockage intermédiaire.

Les ateliers de préparation de minerai du site sont caractérisés par des machines tournantes à haute inertie (concasseurs giratoires, cribles vibrants à double étage, bandes transporteuses motorisées de plusieurs centaines de mètres) soumises à de sévères contraintes mécaniques, vibratoires et thermiques. Ces installations fonctionnent 24h/24 en flux continu, où la moindre défaillance non anticipée perturbe l'ensemble de la chaîne logistique et chimique.

\section{Entité d'Accueil}
Le projet a été hébergé au sein du \textbf{Département Transformation Digitale \& Opérations Industrielles}. Cette entité a pour mission principale :
\begin{itemize}
    \item L'accélération de l'adoption des technologies de l'\textbf{Industrie 4.0} au sein des ateliers de production ;
    \item La conception et l'expérimentation de solutions logicielles souveraines (IIoT, Big Data, Jumeaux Numériques, Intelligence Artificielle) améliorant l'efficacité opérationnelle ;
    \item La réduction des arrêts imprévus d'usines (\textit{Unplanned Downtime}) à travers le déploiement d'outils d'aide à la décision et de maintenance prédictive ;
    \item La valorisation des flux massifs de données générés par les capteurs industriels déployés sur les équipements critiques.
\end{itemize}

\section{Contexte du Stage}
Ce stage d'ingénieur s'inscrit dans le cadre de la 4\textsuperscript{ème} année du cycle d'ingénieur d'État en Génie Informatique à l'\textbf{École Nationale des Sciences Appliquées de Safi (ENSA Safi)}.

L'objectif pédagogique et professionnel est de confronter l'élève ingénieur aux défis réels d'ingénierie logicielle au sein d'un grand groupe industriel, en alliant la théorie des architectures logicielles distribuées (Clean Architecture, DDD, CQRS), l'algorithmique avancée (théorie des graphes) et les technologies Web temps réel (ASP.NET Core 9, WebSockets SignalR, Angular 19).

Le projet confié consistait à concevoir et développer de bout en bout un prototype de \textbf{Jumeau Numérique} capable d'ingérer la télémétrie capteurs, d'identifier instantanément les dérives opérationnelles et de calculer algorithmiquement l'impact d'une panne mécanique sur les machines avals, comblant ainsi le fossé entre la supervision classique et l'aide à la décision intelligente.